% ch04_u02 —— 原书 p.132–p.141（4.1 末、4.2.1–4.2.11）
% 首句承接 ch04_u01 的末句（p.131 末，ch04_u01 以「由 $X$ 上配备……的 $G$-丛 $\mathfrak{F}$ 构成的栈」收尾）；
% 本单元以括注补出 p.132 的省略数据，再接谓语，使拼装后的句子通顺。
\SourcePage{132}
（这里「配备」指的是同构 $\mathcal{F}^{\mathrm{top}}\xrightarrow{\ \sim\ }\xi$ 的同痕类），它非空、连通且单连通的。若把 $X$ 上的 $G$-丛解释为 $X^{\mathrm{top}}$ 上带 $\bar\partial$-联络的 $G$-丛，这一点就很清楚。

\emph{注.} 在 4.1.2 中，我们定义了 $H^{1}(X\setminus S,\Pi(1))$-torsor（主齐性空间）$\phi_{S}\to\operatorname{Bun}_{G}$。若对某个 $x\in X$ 有 $S=\{x\}$，则 $H^{1}(X\setminus S,\Pi(1))=H^{1}(X,\Pi(1))$。故 $\phi_{\{x\}}\to\operatorname{Bun}_{G}$ 是一个 $H^{1}(X,\Pi(1))$-torsor。命题 4.1.11 可以改述如下：若 $\tilde{G}$ 是 $G$ 的泛覆叠，则 Chern 类映射 $\pi_{0}(\operatorname{Bun}_{G})\to\Pi$ 是双射，而 $\phi_{\{x\}}\to\operatorname{Bun}_{G}$ 限制到每个连通分支上是泛覆叠。这确实是改述，因为 $x$ 的选取定义了一个等价。

\begin{equation}\tag{163}
\Pi(1)\text{ gerbes}(X)\xrightarrow{\ \sim\ }\Pi\times H^{1}(X,\Pi(1))\text{ tors}
\end{equation}
（对 $X$ 上的一个 $\Pi(1)$-gerbe，对应它在 $H^{2}(X,\Pi(1))=\Pi$ 中的类，以及它的对象在 $X\setminus\{x\}$ 上的同构类所构成的 $H^{1}(X,\Pi(1))$-torsor）。

\bdsec{4.2. Pfaffian（一）}

本小节我们假定：对 $(\mathbb{Z}/2\mathbb{Z})$-分次向量空间 $A$ 与 $B$，$A\otimes B$ 与 $B\otimes A$ 的等同由 $a\otimes b\mapsto(-1)^{p(a)p(b)}b\otimes a$ 定义，其中 $p(a)$ 是 $a$ 的奇偶性。仿照 [Kn-Mu]，对维数 $n<\infty$ 的向量空间 $V$，我们把 $\det V$ 看作次数为 $n\bmod 2$ 的 $(\mathbb{Z}/2\mathbb{Z})$-分次空间。

\bdp{4.2.1} 设 $X$ 是 $\mathbb{C}$ 上的光滑完备曲线。$X$ 上的 \emph{$\omega$-正交丛}是指带非退化对称配对 $\mathcal{Q}\otimes\mathcal{Q}\to\omega_{X}$ 的向量丛 $\mathcal{Q}$。记 $\omega\text{-Ort}$ 为 $X$ 上 $\omega$-正交丛的栈。在 $\omega\text{-Ort}$ 上有一个熟知的线丛 $\det R\Gamma$（它在 $\mathcal{Q}$ 上的纤维是 $\det R\Gamma(X,\mathcal{Q})$）。Laszlo 与 Sorger [La-So] 在 $\omega\text{-Ort}$ 上构造了一个 $(\mathbb{Z}/2\mathbb{Z})$-分次线丛（他们称之为 \emph{Pfaffian}），并证明 Pfaffian 的张量平方是 $\det R\Gamma$。为我们的目的，使用 Pfaffian 的另一个定义更方便。它当然应该与 [La-So] 中的定义等价，但我们没有验证这一点。

\SourcePage{133}
我们将构造 $\omega\text{-Ort}$ 上的一个线丛 $\operatorname{Pf}$，称之为 \emph{Pfaffian}。它在 $\omega$-正交丛 $\mathcal{Q}$ 上的纤维记作 $\operatorname{Pf}(\mathcal{Q})$。$-1\in\operatorname{Aut}\mathcal{Q}$ 在 $\operatorname{Pf}(\mathcal{Q})$ 上的作用定义了 $\operatorname{Pf}$ 上的一个 $(\mathbb{Z}/2\mathbb{Z})$-分次。由于 $\operatorname{Pf}$ 是线丛，「分次」只意味着存在局部常值函数 $p:(\omega\text{-Ort})\to\mathbb{Z}/2\mathbb{Z}$，使得 $\operatorname{Pf}(\mathcal{Q})$ 的次数为 $p(\mathcal{Q})$。实际上 $p(\mathcal{Q})=\dim H^{0}(\mathcal{Q})\bmod 2$。$\dim H^{0}(\mathcal{Q})\bmod 2$ 局部常值这一事实由 M. Atiyah 与 D. Mumford [At, Mu] 证明。

对 $\omega$-正交丛 $\mathcal{Q}$，记 $\mathcal{Q}^{-}$ 为同一个丛 $\mathcal{Q}$ 配备相反的配对 $\mathcal{Q}\otimes\mathcal{Q}\to\omega_{X}$。置 $\operatorname{Pf}^{-}(\mathcal{Q}):=\operatorname{Pf}(\mathcal{Q}^{-})$。我们将定义一个典范同构 $\operatorname{Pf}\otimes\operatorname{Pf}^{-}\xrightarrow{\ \sim\ }\det R\Gamma$。定义同构 $f_{\pm i}:\operatorname{Pf}(\mathcal{Q})\xrightarrow{\ \sim\ }\operatorname{Pf}(\mathcal{Q}^{-})$ 为 $f_{\pm i}:=(\phi_{\pm i})_{*}$，其中 $i=\sqrt{-1}$，而 $\phi_{i}:\mathcal{Q}\xrightarrow{\ \sim\ }\mathcal{Q}^{-}$ 是乘以 $i$。借助 $f_{\pm i}$ 把 $\operatorname{Pf}$ 与 $\operatorname{Pf}^{-}$ 等同起来，我们得到同构 $c_{\pm i}:\operatorname{Pf}^{\otimes 2}\xrightarrow{\ \sim\ }\det R\Gamma$，使得 $(c_{i})^{-1}c_{-i}:\operatorname{Pf}(\mathcal{Q})^{\otimes 2}\xrightarrow{\ \sim\ }\operatorname{Pf}(\mathcal{Q})^{\otimes 2}$ 是乘以 $(-1)^{p(\mathcal{Q})}$。

\emph{注.}
\begin{enumerate}[label=(\roman*),leftmargin=2.2em]
\item 若 $\mathcal{Q}$ 是 $\omega$-正交丛，则由 Serre 对偶 $H^{1}(X,\mathcal{Q})=(H^{0}(X,\mathcal{Q}))^{*}$，故 $\det R\Gamma(X,\mathcal{Q})=\det H^{0}(X,\mathcal{Q})^{\otimes 2}$。朴素的定义会是 $\operatorname{Pf}^{?}(\mathcal{Q}):=\det H^{0}(X,\mathcal{Q})$，但它对 $\mathcal{Q}$ 的族没有意义，因为 $\dim H^{0}(X,\mathcal{Q})$ 可能跳跃。
\item 设 $\mathcal{Q}$ 是 $\mathcal{Q}_{1}$ 与 $\mathcal{Q}_{2}$ 的正交直和。则 $\det R\Gamma(X,\mathcal{Q})=\det R\Gamma(X,\mathcal{Q}_{1})\otimes\det R\Gamma(X,\mathcal{Q}_{2})$。由 $\operatorname{Pf}$ 与 $\operatorname{Pf}\otimes\operatorname{Pf}^{-}\xrightarrow{\sim}\det R\Gamma$ 的定义将会清楚：存在典范同构 $\operatorname{Pf}(\mathcal{Q})\xrightarrow{\ \sim\ }\operatorname{Pf}(\mathcal{Q}_{1})\otimes\operatorname{Pf}(\mathcal{Q}_{2})$，并且图表
\[
\begin{tikzcd}[column sep=3.2em, row sep=2.6em]
  \operatorname{Pf}(\mathcal{Q})\otimes\operatorname{Pf}(\mathcal{Q}^{-})
    \arrow[r, "\sim"]
    \arrow[d, "\wr"']
  & \operatorname{Pf}(\mathcal{Q}_{1})\otimes\operatorname{Pf}(\mathcal{Q}_{1}^{-})\otimes\operatorname{Pf}(\mathcal{Q}_{2})\otimes\operatorname{Pf}(\mathcal{Q}_{2}^{-})
    \arrow[d, "\wr"] \\
  \det R\Gamma(X,\mathcal{Q})
    \arrow[r, "\sim"]
  & \det R\Gamma(X,\mathcal{Q}_{1})\otimes\det R\Gamma(X,\mathcal{Q}_{2})
\end{tikzcd}
\]
交换。因此同构 $c_{\pm i}:\operatorname{Pf}(\mathcal{Q})^{\otimes 2}\xrightarrow{\ \sim\ }\det R\Gamma(X,\mathcal{Q})$ 与分解 $\mathcal{Q}=\mathcal{Q}_{1}\oplus\mathcal{Q}_{2}$ 相容。
\SourcePage{134}
\item 可以定义 $c_{\pm}:\operatorname{Pf}(\mathcal{Q})^{\otimes 2}\xrightarrow{\ \sim\ }\det R\Gamma(X,\mathcal{Q})$ 为 $c_{\pm}=i^{\pm p(\mathcal{Q})^{2}}c_{i}$，其中 $p(\mathcal{Q})^{2}$ 看作 $\mathbb{Z}/4\mathbb{Z}$ 的元素。于是把 $i$ 换成 $-i$ 时 $c_{\pm}$ 不变。然而 $c_{\pm}$ 似乎不是自然的对象，例如它们与分解 $\mathcal{Q}=\mathcal{Q}_{1}\oplus\mathcal{Q}_{2}$ 不相容（「误差」是 $(-1)^{p(\mathcal{Q}_{1})p(\mathcal{Q}_{2})}$）。
\item 若把 $\mathbb{C}$ 换成满足 $\operatorname{char}k\ne 2$ 的任意域 $k$，$\operatorname{Pf}(\mathcal{Q})$ 的构造仍然有效。$\operatorname{char}k=2$ 的情形在 4.2.16 中讨论。
\end{enumerate}

\bdp{4.2.2} \emph{Lagrange 三元组}由以下数据构成：偶维向量空间 $V$，$V$ 上的非退化双线性对称型 $(\ ,\ )$，以及 Lagrange（= 极大各向同性）子空间 $L_{+},L_{-}\subset V$。若 $X$ 与 $\mathcal{Q}$ 如 4.2.1，而 $\mathcal{Q}'\subset\mathcal{Q}$ 是满足 $H^{0}(X,\mathcal{Q}')=0$ 且 $S:=\operatorname{Supp}(\mathcal{Q}/\mathcal{Q}')$ 有限的子层，则对 $(\mathcal{Q},\mathcal{Q}')$ 如下对应一个 Lagrange 三元组 $(V;L_{+},L_{-})$（参 [Mu]）：
\begin{enumerate}[label=(\arabic*),leftmargin=2.2em]
\item $V:=H^{0}(X,\mathcal{Q}''/\mathcal{Q}')$，其中 $\mathcal{Q}'':=\operatorname{Hom}(\mathcal{Q}',\omega_{X})\supset\mathcal{Q}$；
\item $L_{+}:=H^{0}(X,\mathcal{Q}/\mathcal{Q}')\subset V$；
\item $L_{-}:=H^{0}(X,\mathcal{Q}'')\subset V$；
\item $V$ 上的双线性型由下述两者诱导：自然配对 $\mathcal{Q}''/\mathcal{Q}'\otimes\mathcal{Q}''/\mathcal{Q}'\to(j_{*}\omega_{X\setminus S})/\omega_{X}$，以及「留数之和」映射 $H^{0}(X,(j_{*}\omega_{X\setminus S})/\omega_{X})\to\mathbb{C}$。这里 $j$ 是嵌入 $X\setminus S\to X$。在此情形下，可以把 $R\Gamma(X,\mathcal{Q})$ 与复形
\begin{equation}\tag{164}
0\longrightarrow L_{-}\longrightarrow V/L_{+}\longrightarrow 0
\end{equation}
等同起来，该复形集中位于次数 0 与 1。特别地 $H^{0}(X,\mathcal{Q})=L_{+}\cap L_{-}$，$H^{1}(X,\mathcal{Q})=V/(L_{+}+L_{-})$。而 $H^{0}(X,\mathcal{Q})=L_{+}\cap L_{-}$ 与 $H^{1}(X,\mathcal{Q})=V/(L_{+}+L_{-})$ 之间的 Serre 配对由 $V$ 上的双线性型诱导。
\end{enumerate}

\bdp{4.2.3} 对 Lagrange 三元组 $(V;L_{+},L_{-})$，置
\begin{equation}\tag{165}
\det(V;L_{+},L_{-}):=\det L_{+}\otimes\det L_{-}\otimes(\det V)^{*}.
\end{equation}

\SourcePage{135}
$\det(V;L_{+},L_{-})$ 正是复形 (164) 的行列式。公式 (165) 在 Lagrange 三元组的栈上定义了一个线丛 $\det$。在 4.2.4 与 4.2.8 中，我们将在此栈上构造一个 $(\mathbb{Z}/2\mathbb{Z})$-分次线丛 $\operatorname{Pf}$，以及一个典范同构 $\operatorname{Pf}\otimes\operatorname{Pf}^{-}\xrightarrow{\ \sim\ }\det$。这里 $\operatorname{Pf}^{-}(V;L_{+},L_{-}):=\operatorname{Pf}(V^{-};L_{+},L_{-})$，而 $V^{-}$ 表示 $V$ 配备型 $-(\ ,\ )$。朴素的「定义」会是 $\operatorname{Pf}^{?}(V;L_{+},L_{-}):=\det(L_{+}\cap L_{-})$，或 $\operatorname{Pf}^{?}(V;L_{+},L_{-})^{*}:=\det((L_{+}\cap L_{-})^{*})=\det(V/(L_{+}+L_{-}))$（参 4.2.1 的注 (i)）。

\bdp{4.2.4} 对 Lagrange 三元组 $(V;L_{+},L_{-})$，如下定义 $\operatorname{Pf}(V;L_{+},L_{-})$。记 $\operatorname{Cl}(V)$ 为带典范 $(\mathbb{Z}/2\mathbb{Z})$-分次的 Clifford 代数（$V\subset\operatorname{Cl}(V)$ 为奇）。设 $M$ 是不可约 $(\mathbb{Z}/2\mathbb{Z})$-分次 $\operatorname{Cl}(V)$-模（实际上，即使不计分次，$M$ 也不可约）。$M$ 在张量一维 $(\mathbb{Z}/2\mathbb{Z})$-分次向量空间的意义下唯一。置 $M_{L_{-}}=M/L_{-}M$，$M^{L_{+}}:=\{m\in M\mid L_{+}m=0\}$。则 $M^{L_{+}}$ 与 $M_{L_{-}}$ 是一维 $(\mathbb{Z}/2\mathbb{Z})$-分次空间。我们置
\begin{equation}\tag{166}
\operatorname{Pf}(V;L_{+},L_{-}):=M^{L_{+}}\otimes(M_{L_{-}})^{*}.
\end{equation}
特别地，可以取 $M=\operatorname{Cl}(V)/\operatorname{Cl}(V)L_{+}$。此时 $M^{L_{+}}=\mathbb{C}$，故
\begin{equation}\tag{167}
\operatorname{Pf}(V;L_{+},L_{-})^{*}=\operatorname{Cl}(V)/(L_{-}\cdot\operatorname{Cl}(V)+\operatorname{Cl}(V)\cdot L_{+}).
\end{equation}
显然 (166) 或 (167) 把 $\operatorname{Pf}$ 定义为 Lagrange 三元组栈上的 $(\mathbb{Z}/2\mathbb{Z})$-分次线丛。\footnote{换言之，从个别 Lagrange 三元组过渡到族是显然的。此原则贯穿我们对 Pfaffian 的全部讨论（只有在 4.2.14 的无穷维情形中，我们才显式考虑族，因为这确实需要一些小心）。} 该分次对应于 $-1\in\operatorname{Aut}(V;L_{+},L_{-})$ 在 $\operatorname{Pf}(V;L_{+},L_{-})$ 上的作用。

若 $V$ 是 $V_{1}$ 与 $V_{2}$ 的正交直和，则 $\operatorname{Cl}(V)$ 是超代数 $\operatorname{Cl}(V_{1})$ 与 $\operatorname{Cl}(V_{2})$ 的张量积。因此，若 $(V^{1};L^{1}_{+},L^{1}_{-})$ 与 $(V^{2};L^{2}_{+},L^{2}_{-})$ 是 Lagrange 三元组，则有典范同构
\begin{equation}\tag{168}
\begin{aligned}
&\operatorname{Pf}(V^{1}\oplus V^{2};L^{1}_{+}\oplus L^{2}_{+},L^{1}_{-}\oplus L^{2}_{-})\\
&\qquad=\operatorname{Pf}(V^{1};L^{1}_{+},L^{1}_{-})\otimes\operatorname{Pf}(V^{2};L^{2}_{+},L^{2}_{-}),
\end{aligned}
\end{equation}
其中 $\oplus$ 表示正交直和。

\SourcePage{136}
$\operatorname{Pf}(V;L_{+},L_{-})$ 为偶，当且仅当 $\dim(L_{+}\cap L_{-})$ 为偶。这由 (168) 与下面引理的陈述 (i) 推出。

\bdp{4.2.5} \emph{引理.}
\begin{enumerate}[label=(\roman*),leftmargin=2.2em]
\item 任意 Lagrange 三元组 $(V;L_{+},L_{-})$ 都可表为 Lagrange 三元组 $(V^{1};L^{1}_{+},L^{1}_{-})$ 与 $(V^{2};L^{2}_{+},L^{2}_{-})$ 的正交直和，使得 $L^{1}_{+}\cap L^{1}_{-}=0$，$L^{2}_{+}=L^{2}_{-}$。
\item 此外，若固定子空间 $\Lambda\subset L_{+}$ 使 $L_{+}=\Lambda\oplus(L_{+}\cap L_{-})$，则可选取上述分解 $(V;L_{+},L_{-})=(V^{1};L^{1}_{+},L^{1}_{-})\oplus(V^{2};L^{2}_{+},L^{2}_{-})$，使得 $L^{1}_{+}=\Lambda$。
\end{enumerate}

\emph{证明.}
\begin{enumerate}[label=(\roman*),leftmargin=2.2em]
\item 选取子空间 $P\subset V$ 使 $V=(L_{+}+L_{-})\oplus P$。然后置 $V^{2}:=(L_{1}\cap L_{2})\oplus P$，$V^{1}:=(V^{2})^{\perp}$。
% [Tutor Note: 原书 p.136 证明 (i) 中印作 $L_1\cap L_2$，按上下文应为 $L_+\cap L_-$；底本如此，保留原样。]
\item 选取子空间 $P\subset\Lambda^{\perp}$ 使 $\Lambda^{\perp}=L_{+}\oplus P$（这意味着 $V=(L_{+}+L_{-})\oplus P$，因为 $\Lambda^{\perp}/L_{+}\to V/(L_{+}+L_{-})$ 是同构）。然后如上进行。\hfill$\Box$
\end{enumerate}

\bdp{4.2.6} 本小节（读者可以跳过）我们在 $\operatorname{Pf}(V;L_{+},L_{-})$ 与 4.2.3 的朴素 $\operatorname{Pf}^{?}(V;L_{+},L_{-})$ 之间构造一个典范同构。回忆 $\operatorname{Pf}^{?}(V;L_{+},L_{-}):=\det(L_{+}\cap L_{-})$，故 $\operatorname{Pf}^{?}(V;L_{+},L_{-})^{*}=\det((L_{+}\cap L_{-})^{*})=\det(V/(L_{+}+L_{-}))$。这里应理解为：配对 $\det W\otimes\det W^{*}\to\mathbb{C}$（$W:=L_{+}\cap L_{-}$）由 $(e_{1}\wedge\dots\wedge e_{k})\otimes(e^{k}\wedge\dots\wedge e^{1})\mapsto 1$ 定义，其中 $e_{1},\dots,e_{k}$ 是 $W$ 的基，$e^{1},\dots,e^{k}$ 是 $W^{*}$ 的对偶基（从「超」的观点看，这个配对是合理的；例如它与分解 $W=W_{1}\oplus W_{2}$ 相容）。

为定义同构 $\operatorname{Pf}(V;L_{+},L_{-})\xrightarrow{\ \sim\ }\operatorname{Pf}^{?}(V;L_{+},L_{-})$，我们使用 $\operatorname{Cl}(V)$ 上由
\begin{equation}\tag{169}
\operatorname{Cl}_{0}(V)=\mathbb{C},\qquad \operatorname{Cl}_{k+1}(V)=\operatorname{Cl}_{k}(V)+V\cdot\operatorname{Cl}_{k}(V)
\end{equation}
定义的典范滤链。我们有 $\operatorname{Cl}_{k}(V)/\operatorname{Cl}_{k-1}(V)=\bigwedge^{k}V$。置 $r:=\dim(L_{+}\cap L_{-})$。有典范满射
% 下式起于 p.136 末、止于 p.137 首，跨页处只拆行不拆义。
\SourcePage{137}
\[
\begin{aligned}
\varphi:\operatorname{Cl}_{r}(V)&\longrightarrow\bigwedge\nolimits^{r}V\longrightarrow\bigwedge\nolimits^{r}(V/(L_{+}+L_{-}))\\
&=\det(V/(L_{+}+L_{-}))=\operatorname{Pf}^{?}(V;L_{+},L_{-})^{*}.
\end{aligned}
\]
由 4.2.5(i) 容易推出：典范映射 $\operatorname{Cl}_{r}(V)\to\operatorname{Cl}(V)/(L_{-}\cdot\operatorname{Cl}(V)+\operatorname{Cl}(V)\cdot L_{+})=\operatorname{Pf}(V;L_{+},L_{-})^{*}$ 经过 $\varphi$ 分解。由此诱导的映射 $f:\operatorname{Pf}^{?}(V;L_{+},L_{-})^{*}\to\operatorname{Pf}(V;L_{+},L_{-})^{*}$ 是同构。$f^{*}$ 就是所求的同构 $\operatorname{Pf}(V;L_{+},L_{-})\xrightarrow{\ \sim\ }\operatorname{Pf}^{?}(V;L_{+},L_{-})$。

这里有一个等价定义。设 $M$ 是不可约 $(\mathbb{Z}/2\mathbb{Z})$-分次 $\operatorname{Cl}(V)$-模。典范嵌入 $\det(L_{+}\cap L_{-})\subset\bigwedge^{*}(L_{+}\cap L_{-})=\operatorname{Cl}(L_{+}\cap L_{-})\subset\operatorname{Cl}(V)$ 诱导映射 $\det(L_{+}\cap L_{-})\otimes M_{L_{+}\cap L_{-}}\to M^{L_{+}\cap L_{-}}$，它实际上是同构。容易由 4.2.5(i) 推出，复合
\[
\begin{aligned}
M^{L_{+}}\longrightarrow M^{L_{+}\cap L_{-}}&\xrightarrow{\ \sim\ }\det(L_{+}\cap L_{-})\otimes M_{L_{+}\cap L_{-}}\\
&\longrightarrow\det(L_{+}\cap L_{-})\otimes M_{L_{-}}
\end{aligned}
\]
是同构。它诱导同构 $\operatorname{Pf}(V;L_{+},L_{-}):=M^{L_{+}}\otimes(M_{L_{-}})^{\otimes-1}\to\det(L_{+}\cap L_{-})=\operatorname{Pf}^{?}(V;L_{+},L_{-})$，这实际上与上面构造的同构互逆。

\bdp{4.2.7} 在构造同构 $\operatorname{Pf}\otimes\operatorname{Pf}^{-}\xrightarrow{\ \sim\ }\det$ 之前，我们先构造一个典范同构
\begin{equation}\tag{170}
\begin{aligned}
&\operatorname{Pf}(V\oplus V^{*};L_{+}\oplus L_{+}^{\perp},L_{-}\oplus L_{-}^{\perp})\\
&\qquad\xrightarrow{\ \sim\ }\det(V;L_{+},L_{-}),
\end{aligned}
\end{equation}
其中 $V$ 是不带任何双线性型的有限维向量空间，$L_{\pm}\subset V$ 是任意子空间，而 $V\oplus V^{*}$ 配备显然的双线性型（(170) 的左端有意义，因为 $L_{\pm}\oplus L_{\pm}^{\perp}$ 是 Lagrange 的；(170) 的右端由 (165) 定义）。置
\begin{equation}\tag{171}
M=\bigwedge V\otimes(\det L_{+})^{*},\qquad \bigwedge V:=\bigoplus_{i}\bigwedge\nolimits^{i}V.
\end{equation}
$M$ 是不可约 $\operatorname{Cl}(V\oplus V^{*})$-模，且 $M^{L_{+}\oplus L_{+}^{\perp}}=\mathbb{C}$，故按 (166)，$\operatorname{Pf}(V\oplus V^{*};L_{+}\oplus L_{+}^{\perp},L_{-}\oplus L_{-}^{\perp})=(M_{L_{-}\oplus L_{-}^{\perp}})^{*}$。显然 $M_{L_{-}}=\bigwedge(V/L_{-})\otimes(\det L_{+})^{*}$，且 $M_{L_{-}\oplus L_{-}^{\perp}}=\det(V/L_{-})\otimes(\det L_{+})^{*}=\det(V;L_{+},L_{-})^{*}$（见 (165)）。于是我们构造了同构 (170)。

\SourcePage{138}
\bdp{4.2.8} 现在设 $(V;L_{+},L_{-})$ 是 Lagrange 三元组。我们将构造一个典范同构
\begin{equation}\tag{172}
\begin{aligned}
&\operatorname{Pf}(V;L_{+},L_{-})\otimes\operatorname{Pf}(V^{-};L_{+},L_{-})\\
&\qquad\xrightarrow{\ \sim\ }\det(V;L_{+},L_{-}),
\end{aligned}
\end{equation}
其中 $V^{-}$ 表示 $V$ 配备双线性型 $-(\ ,\ )$。若 $W$ 是配备非退化对称双线性型的有限维向量空间，则 $(V\otimes W;L_{+}\otimes W,L_{-}\otimes W)$ 是 Lagrange 三元组。(170) 可以改写为典范同构
\begin{equation}\tag{173}
\det(V;L_{+},L_{-})\xrightarrow{\ \sim\ }\operatorname{Pf}(V\otimes H;L_{+}\otimes H,L_{-}\otimes H),
\end{equation}
其中 $H$ 表示 $\mathbb{C}^{2}$ 配备双线性型 $\begin{pmatrix}0&1\\1&0\end{pmatrix}$。另一方面，(168) 给出同构
\begin{equation}\tag{174}
\begin{aligned}
&\operatorname{Pf}(V;L_{+},L_{-})\otimes\operatorname{Pf}(V^{-};L_{+},L_{-})\\
&\qquad\xrightarrow{\ \sim\ }\operatorname{Pf}(V\otimes H';L_{+}\otimes H',L_{-}\otimes H'),
\end{aligned}
\end{equation}
其中 $H'$ 表示 $\mathbb{C}^{2}$ 配备双线性型 $\begin{pmatrix}1&0\\0&-1\end{pmatrix}$。于是同构 $\varphi:H'\xrightarrow{\ \sim\ }H$ 诱导同构
\[
\varphi_{*}:\operatorname{Pf}(V;L_{+},L_{-})\otimes\operatorname{Pf}(V^{-};L_{+},L_{-})\xrightarrow{\ \sim\ }\det(V;L_{+},L_{-}).
\]
\emph{引理.} 若 $\psi\in\operatorname{Aut}H'$，则
\begin{equation}\tag{175}
(\varphi\psi)_{*}=(\det\psi)^{n}\varphi_{*},\qquad n=\dim(L_{+}\cap L_{-}).
\end{equation}
\emph{证明.} $\operatorname{Aut}H'$ 通过某个特征标 $\chi:\operatorname{Aut}H'\to\mathbb{C}^{*}$ 作用在 (174) 的右端。$\operatorname{Aut}H'$ 的任意特征标都具有形式 $\psi\mapsto(\det\psi)^{m}$，$m\in\mathbb{Z}/2\mathbb{Z}$。$\chi\begin{pmatrix}1&0\\0&-1\end{pmatrix}=(-1)^{n}$，$n:=\dim(L_{+}\cap L_{-})$，因为 $-1\in\operatorname{Aut}(V;L_{+},L_{-})$ 在 $\operatorname{Pf}(V;L_{+},L_{-})$ 上的作用是 $(-1)^{n}$（见 4.2.4）。故 $m=n\bmod 2$。\hfill$\Box$

我们如下定义 (172)：取任意满足 $\det\varphi=1$ 的同构 $\varphi:H'\xrightarrow{\ \sim\ }H$，令 (172) 为 $\varphi_{*}$。

\emph{注.}
\begin{enumerate}[label=(\roman*),leftmargin=2.2em]
\item (172) 与 $(V;L_{+},L_{-})$ 到正交直和的分解相容；也就是说，若有一个这样的分解 $(V;L_{+},L_{-})=(V^{1};L^{1}_{+},L^{1}_{-})\oplus(V^{2};L^{2}_{+},L^{2}_{-})$，则同构\SourcePage{139}
(172) 对 $(V;L_{+},L_{-})$、$(V^{1};L^{1}_{+},L^{1}_{-})$ 与 $(V^{2};L^{2}_{+},L^{2}_{-})$ 这三个三元组，都与 (168) 及典范同构 $\det(V;L_{+},L_{-})=\det(V^{1};L^{1}_{+},L^{1}_{-})\otimes\det(V^{2};L^{2}_{+},L^{2}_{-})$ 相容。
\item (170) 与 $(V;L_{+},L_{-})$ 到直和的分解相容。
\end{enumerate}

\bdp{4.2.9} 本小节（读者可以跳过）给出 (172) 的一个等价构造。我们将使用超代数反同构 $*:\operatorname{Cl}(V^{-})\xrightarrow{\ \sim\ }\operatorname{Cl}(V)$，它在 $V$ 上恒等（对任意 $v_{1},\dots,v_{k}\in V$，有 $(v_{1}\cdots v_{k})^{*}=(-1)^{k(k-1)/2}v_{k}\cdots v_{1}$）。我们还使用典范映射 $\operatorname{sTr}:\operatorname{Cl}(V)=\operatorname{Cl}_{n}(V)\to\operatorname{Cl}_{n}(V)/\operatorname{Cl}_{n-1}(V)=\det V$，其中 $n=\dim V$，而 $\operatorname{Cl}_{k}(V)$ 由 (169) 定义。它具有「超迹性质」
\begin{equation}\tag{176}
\operatorname{sTr}(ab)=(-1)^{p(a)p(b)}\operatorname{sTr}(ba),
\end{equation}
其中 $a,b\in\operatorname{Cl}(V)$ 是次数为 $p(a),p(b)\in\mathbb{Z}/2\mathbb{Z}$ 的齐次元素。事实上，只需在 $a\in V$、$p(ab)=n\bmod 2$ 的情形证明 (176)。此时 $b\in\operatorname{Cl}_{n-1}(V)$，(176) 显然。或者可以验证：$\operatorname{sTr}(a)$ 与算子 $a:M\to M$ 的超迹至多相差一个符号，其中 $M$ 是不可约 $\operatorname{Cl}(V)$-模。

现在考虑映射
\begin{equation}\tag{177}
\begin{aligned}
&\det L_{-}\otimes\operatorname{Pf}(V;L_{+},L_{-})^{*}\otimes\det L_{+}\otimes\operatorname{Pf}(V^{-};L_{+},L_{-})^{*}\\
&\qquad\longrightarrow\det V,
\end{aligned}
\end{equation}
它由 $a_{-}\otimes x\otimes a_{+}\otimes y\mapsto\operatorname{sTr}(a_{-}xa_{+}y^{*})$ 定义。这里 $a_{\pm}\in\det L_{\pm}\subset\bigwedge^{*}(L_{\pm})=\operatorname{Cl}(L_{\pm})\subset\operatorname{Cl}(V)$，$x\in\operatorname{Pf}(V;L_{+},L_{-})^{*}=\operatorname{Cl}(V)/(L_{-}\cdot\operatorname{Cl}(V)+\operatorname{Cl}(V)\cdot L_{+})$，$y^{*}\in\operatorname{Cl}(V)/(L_{+}\cdot\operatorname{Cl}(V)+\operatorname{Cl}(V)\cdot L_{-})$，故 (177) 是良定义的。容易看出（例如由 4.2.5(i)），(177) 是同构。它诱导同构
\[
\operatorname{Pf}(V;L_{+},L_{-})\otimes\operatorname{Pf}(V^{-};L_{+},L_{-})\xrightarrow{\ \sim\ }\det L_{+}\otimes\det L_{-}\otimes(\det V)^{*}=\det(V;L_{+},L_{-}).
\]
可以证明，这个同构等于 (172)。

\SourcePage{140}
\bdp{4.2.10} 设 $X$ 与 $\mathcal{Q}$ 如 4.2.1，$\mathcal{Q}'\subset\mathcal{Q}$ 如 4.2.2。对这些数据，我们已对应一个 Lagrange 三元组 $(V;L_{+},L_{-})$，使得 $\det(V;L_{+},L_{-})=\det R\Gamma(X,\mathcal{Q})$（见 4.2.2）。置 $\operatorname{Pf}_{\mathcal{Q}'}(\mathcal{Q}):=\operatorname{Pf}(V;L_{+},L_{-})$。按 4.2.9，有典范同构 $\operatorname{Pf}_{\mathcal{Q}'}(\mathcal{Q})\otimes\operatorname{Pf}_{\mathcal{Q}'}(\mathcal{Q}^{-})\xrightarrow{\ \sim\ }\det R\Gamma(X,\mathcal{Q})$。为定义 $\operatorname{Pf}(\mathcal{Q})$，只需对所有满足 $\mathcal{Q}'\subset\tilde{\mathcal{Q}}'$ 的对 $(\mathcal{Q}',\tilde{\mathcal{Q}}')$，定义相容的同构系统 $\operatorname{Pf}_{\mathcal{Q}'}(\mathcal{Q})\xrightarrow{\ \sim\ }\operatorname{Pf}_{\tilde{\mathcal{Q}}'}(\mathcal{Q})$。为定义 $\operatorname{Pf}(\mathcal{Q})\otimes\operatorname{Pf}(\mathcal{Q}^{-})\xrightarrow{\ \sim\ }\det R\Gamma(X,\mathcal{Q})$，只需证明下述图表的交换性：
% [Tutor Note: 原书 p.140 此交换图排印错乱：底行右端的对象缺失，且两支之间的纵向箭头被印成横箭头（落在顶行张量积号下方）。译文把该箭头排为纵向，底行右端照原样留空。]
\[
\begin{tikzcd}[column sep=3.4em, row sep=2.6em]
  \operatorname{Pf}_{\mathcal{Q}'}(\mathcal{Q})\otimes\operatorname{Pf}_{\mathcal{Q}'}(\mathcal{Q}^{-})
    \arrow[r, "\sim"]
    \arrow[d]
  & \det R\Gamma(X,\mathcal{Q}) \\
  \operatorname{Pf}_{\tilde{\mathcal{Q}}'}(\mathcal{Q})\otimes\operatorname{Pf}_{\tilde{\mathcal{Q}}'}(\mathcal{Q}^{-})
    \arrow[r, "\sim"]
  & \phantom{\det R\Gamma(X,\mathcal{Q})}
\end{tikzcd}
\]
对应 $\tilde{\mathcal{Q}}'$ 的 Lagrange 三元组 $(\tilde{V};\tilde{L}_{+},\tilde{L}_{-})$ 与对应 $\mathcal{Q}'$ 的三元组 $(V;L_{+},L_{-})$ 如下相关：若 $\Lambda=H^{0}(X,\tilde{\mathcal{Q}}'/\mathcal{Q}')\subset H^{0}(X,\mathcal{Q}/\mathcal{Q}')=L_{+}$，则
\begin{equation}\tag{178}
\tilde{V}=\Lambda^{\perp}/\Lambda,\qquad \tilde{L}_{+}=L_{+}/\Lambda\subset\tilde{V},\qquad \tilde{L}_{-}=L_{-}\cap\Lambda^{\perp}\hookrightarrow\tilde{V}
\end{equation}
（注意 $\Lambda\cap L_{-}=H^{0}(X,\tilde{\mathcal{Q}}')=0$）。于是剩下的是一些线性代数（见 4.2.11）。容易验证，我们对 $\operatorname{Pf}(\mathcal{Q})$ 以及 $\operatorname{Pf}(\mathcal{Q})\otimes\operatorname{Pf}(\mathcal{Q}^{-})\xrightarrow{\ \sim\ }\det R\Gamma(X,\mathcal{Q})$ 的定义对 $\mathcal{Q}$ 的族有意义。

\bdp{4.2.11} 设 $(V;L_{+},L_{-})$ 是 Lagrange 三元组，$\Lambda\subset L_{+}$ 是满足 $\Lambda\cap L_{-}=0$ 的子空间。则由 (178) 定义的 $(\tilde{V};\tilde{L}_{+},\tilde{L}_{-})$ 是 Lagrange 三元组。在此情形下，我们说 $(\tilde{V};\tilde{L}_{+},\tilde{L}_{-})$ 是 $(V;L_{+},L_{-})$ 的 \emph{子商}。容易证明，子商的子商仍是子商。于是可以考虑以 Lagrange 三元组为对象的范畴 $T$，其中从 $(V;L_{+},L_{-})$ 到 $(V';L'_{+},L'_{-})$ 的态射定义为 $(V;L_{+},L_{-})$ 与 $(V';L'_{+},L'_{-})$ 的某个子商之间的同构。再考虑范畴 $C$，其对象是有限维向量空间的有限复形，态射是拟同构。记 $\mathbb{I}$ 为如下范畴：对象是 $(\mathbb{Z}/2\mathbb{Z})$-分次一维向量空间，态射是保持\SourcePage{141}%
分次的同构。

复形 (164) 作为 $C$ 的对象，函子性地依赖于 $(V;L_{+},L_{-})\in T$：若 $(\tilde{V};\tilde{L}_{+},\tilde{L}_{-})$ 是对应于 $\Lambda\subset L_{+}$ 的 $(V;L_{+},L_{-})$ 的子商，则有拟同构
% [Tutor Note: 原书 p.141 此图把两支纵向的嵌入箭头印成横向钩箭头（分置两行之间偏左、偏右处），未与上下两行对齐；译文按图意排为纵向 hook 箭头：$\tilde{L}_-\hookrightarrow L_-$ 与 $\Lambda^\perp/L_+\hookrightarrow V/L_+$。]
\[
\begin{tikzcd}[column sep=3.6em, row sep=2.4em]
  L_{-} \arrow[r] & V/L_{+} \\
  \tilde{L}_{-} \arrow[r] \arrow[u, hook] & \tilde{V}/\tilde{L}_{+}=\Lambda^{\perp}/L_{+} \arrow[u, hook]
\end{tikzcd}
\]
应用 [Kn-Mu] 的函子 $\det:C\to\mathbb{I}$，我们看到 $\det(V;L_{+},L_{-})\in\mathbb{I}$ 函子性地依赖于 $(V;L_{+},L_{-})\in T$。若 $(\tilde{V};\tilde{L}_{+},\tilde{L}_{-})$ 是对应于 $\Lambda\subset L_{+}$ 的 $(V;L_{+},L_{-})$ 的子商，则 $\det(V;L_{+},L_{-})=(\det L_{+})\otimes(\det L_{-})\otimes(\det V)^{*}$ 与 $\det(\tilde{V};\tilde{L}_{+},\tilde{L}_{-})=(\det\tilde{L}_{+})\otimes(\det\tilde{L}_{-})\otimes(\det\tilde{V})^{*}$ 之间的同构来自自然同构 $\det L_{+}=\det\Lambda\otimes\det\tilde{L}_{+}$、$\det L_{-}=\det\tilde{L}_{-}\otimes\det(V/\Lambda^{\perp})$ 与 $\det V=\det\Lambda\otimes\det\tilde{V}\otimes\det(V/\Lambda^{\perp})$。

正如 4.2.10 所解释的，我们必须把 $\operatorname{Pf}$ 定义为函子 $T\to\mathbb{I}$，并证明 4.2.8 中的同构 $\operatorname{Pf}(V;L_{+},L_{-})\otimes\operatorname{Pf}^{-}(V;L_{+},L_{-})\xrightarrow{\ \sim\ }\det(V;L_{+},L_{-})$ 具有函子性。

若 $(\tilde{V};\tilde{L}_{+},\tilde{L}_{-})$ 是对应于 $\Lambda\subset L_{+}$ 的 $(V;L_{+},L_{-})$ 的子商，则
\[
\begin{aligned}
\operatorname{Pf}(V;L_{+},L_{-})^{*}&=\operatorname{Cl}(V)/(L_{-}\cdot\operatorname{Cl}(V)+\operatorname{Cl}(V)\cdot L_{+}),\\
\operatorname{Pf}(\tilde{V};\tilde{L}_{+},\tilde{L}_{-})^{*}
&=\operatorname{Cl}(\Lambda^{\perp})/\bigl((L_{-}\cap\Lambda^{\perp})\cdot\operatorname{Cl}(\Lambda^{\perp})\\
&\qquad+\operatorname{Cl}(\Lambda^{\perp})\cdot L_{+}\bigr).
\end{aligned}
\]
于是嵌入 $\operatorname{Cl}(\Lambda^{\perp})\to\operatorname{Cl}(V)$ 诱导映射
\begin{equation}\tag{179}
\operatorname{Pf}(\tilde{V};\tilde{L}_{+},\tilde{L}_{-})^{*}\to\operatorname{Pf}(V;L_{+},L_{-})^{*}.
\end{equation}
这定义了 $\operatorname{Pf}^{*}$ 为函子 $T\to\{(\mathbb{Z}/2\mathbb{Z})\text{-分次一维空间}\}$（容易看出，复合对应复合）。还需证明
\begin{enumerate}[label=\alph*),leftmargin=2.2em]
\item (179) 是同构，
\end{enumerate}
